% Options for packages loaded elsewhere
\PassOptionsToPackage{unicode}{hyperref}
\PassOptionsToPackage{hyphens}{url}
\documentclass[
]{article}
\usepackage{xcolor}
\usepackage[margin=1in]{geometry}
\usepackage{amsmath,amssymb}
\setcounter{secnumdepth}{5}
\usepackage{iftex}
\ifPDFTeX
  \usepackage[T1]{fontenc}
  \usepackage[utf8]{inputenc}
  \usepackage{textcomp} % provide euro and other symbols
\else % if luatex or xetex
  \usepackage{unicode-math} % this also loads fontspec
  \defaultfontfeatures{Scale=MatchLowercase}
  \defaultfontfeatures[\rmfamily]{Ligatures=TeX,Scale=1}
\fi
\usepackage{lmodern}
\ifPDFTeX\else
  % xetex/luatex font selection
\fi
% Use upquote if available, for straight quotes in verbatim environments
\IfFileExists{upquote.sty}{\usepackage{upquote}}{}
\IfFileExists{microtype.sty}{% use microtype if available
  \usepackage[]{microtype}
  \UseMicrotypeSet[protrusion]{basicmath} % disable protrusion for tt fonts
}{}
\makeatletter
\@ifundefined{KOMAClassName}{% if non-KOMA class
  \IfFileExists{parskip.sty}{%
    \usepackage{parskip}
  }{% else
    \setlength{\parindent}{0pt}
    \setlength{\parskip}{6pt plus 2pt minus 1pt}}
}{% if KOMA class
  \KOMAoptions{parskip=half}}
\makeatother
\usepackage{color}
\usepackage{fancyvrb}
\newcommand{\VerbBar}{|}
\newcommand{\VERB}{\Verb[commandchars=\\\{\}]}
\DefineVerbatimEnvironment{Highlighting}{Verbatim}{commandchars=\\\{\}}
% Add ',fontsize=\small' for more characters per line
\usepackage{framed}
\definecolor{shadecolor}{RGB}{248,248,248}
\newenvironment{Shaded}{\begin{snugshade}}{\end{snugshade}}
\newcommand{\AlertTok}[1]{\textcolor[rgb]{0.94,0.16,0.16}{#1}}
\newcommand{\AnnotationTok}[1]{\textcolor[rgb]{0.56,0.35,0.01}{\textbf{\textit{#1}}}}
\newcommand{\AttributeTok}[1]{\textcolor[rgb]{0.13,0.29,0.53}{#1}}
\newcommand{\BaseNTok}[1]{\textcolor[rgb]{0.00,0.00,0.81}{#1}}
\newcommand{\BuiltInTok}[1]{#1}
\newcommand{\CharTok}[1]{\textcolor[rgb]{0.31,0.60,0.02}{#1}}
\newcommand{\CommentTok}[1]{\textcolor[rgb]{0.56,0.35,0.01}{\textit{#1}}}
\newcommand{\CommentVarTok}[1]{\textcolor[rgb]{0.56,0.35,0.01}{\textbf{\textit{#1}}}}
\newcommand{\ConstantTok}[1]{\textcolor[rgb]{0.56,0.35,0.01}{#1}}
\newcommand{\ControlFlowTok}[1]{\textcolor[rgb]{0.13,0.29,0.53}{\textbf{#1}}}
\newcommand{\DataTypeTok}[1]{\textcolor[rgb]{0.13,0.29,0.53}{#1}}
\newcommand{\DecValTok}[1]{\textcolor[rgb]{0.00,0.00,0.81}{#1}}
\newcommand{\DocumentationTok}[1]{\textcolor[rgb]{0.56,0.35,0.01}{\textbf{\textit{#1}}}}
\newcommand{\ErrorTok}[1]{\textcolor[rgb]{0.64,0.00,0.00}{\textbf{#1}}}
\newcommand{\ExtensionTok}[1]{#1}
\newcommand{\FloatTok}[1]{\textcolor[rgb]{0.00,0.00,0.81}{#1}}
\newcommand{\FunctionTok}[1]{\textcolor[rgb]{0.13,0.29,0.53}{\textbf{#1}}}
\newcommand{\ImportTok}[1]{#1}
\newcommand{\InformationTok}[1]{\textcolor[rgb]{0.56,0.35,0.01}{\textbf{\textit{#1}}}}
\newcommand{\KeywordTok}[1]{\textcolor[rgb]{0.13,0.29,0.53}{\textbf{#1}}}
\newcommand{\NormalTok}[1]{#1}
\newcommand{\OperatorTok}[1]{\textcolor[rgb]{0.81,0.36,0.00}{\textbf{#1}}}
\newcommand{\OtherTok}[1]{\textcolor[rgb]{0.56,0.35,0.01}{#1}}
\newcommand{\PreprocessorTok}[1]{\textcolor[rgb]{0.56,0.35,0.01}{\textit{#1}}}
\newcommand{\RegionMarkerTok}[1]{#1}
\newcommand{\SpecialCharTok}[1]{\textcolor[rgb]{0.81,0.36,0.00}{\textbf{#1}}}
\newcommand{\SpecialStringTok}[1]{\textcolor[rgb]{0.31,0.60,0.02}{#1}}
\newcommand{\StringTok}[1]{\textcolor[rgb]{0.31,0.60,0.02}{#1}}
\newcommand{\VariableTok}[1]{\textcolor[rgb]{0.00,0.00,0.00}{#1}}
\newcommand{\VerbatimStringTok}[1]{\textcolor[rgb]{0.31,0.60,0.02}{#1}}
\newcommand{\WarningTok}[1]{\textcolor[rgb]{0.56,0.35,0.01}{\textbf{\textit{#1}}}}
\usepackage{longtable,booktabs,array}
\usepackage{calc} % for calculating minipage widths
% Correct order of tables after \paragraph or \subparagraph
\usepackage{etoolbox}
\makeatletter
\patchcmd\longtable{\par}{\if@noskipsec\mbox{}\fi\par}{}{}
\makeatother
% Allow footnotes in longtable head/foot
\IfFileExists{footnotehyper.sty}{\usepackage{footnotehyper}}{\usepackage{footnote}}
\makesavenoteenv{longtable}
\usepackage{graphicx}
\makeatletter
\newsavebox\pandoc@box
\newcommand*\pandocbounded[1]{% scales image to fit in text height/width
  \sbox\pandoc@box{#1}%
  \Gscale@div\@tempa{\textheight}{\dimexpr\ht\pandoc@box+\dp\pandoc@box\relax}%
  \Gscale@div\@tempb{\linewidth}{\wd\pandoc@box}%
  \ifdim\@tempb\p@<\@tempa\p@\let\@tempa\@tempb\fi% select the smaller of both
  \ifdim\@tempa\p@<\p@\scalebox{\@tempa}{\usebox\pandoc@box}%
  \else\usebox{\pandoc@box}%
  \fi%
}
% Set default figure placement to htbp
\def\fps@figure{htbp}
\makeatother
\setlength{\emergencystretch}{3em} % prevent overfull lines
\providecommand{\tightlist}{%
  \setlength{\itemsep}{0pt}\setlength{\parskip}{0pt}}
\usepackage{bookmark}
\IfFileExists{xurl.sty}{\usepackage{xurl}}{} % add URL line breaks if available
\urlstyle{same}
\hypersetup{
  pdftitle={Lab 01 -- Data Anatomy \& Type Checks},
  pdfauthor={Your Name},
  hidelinks,
  pdfcreator={LaTeX via pandoc}}

\title{Lab 01 -- Data Anatomy \& Type Checks}
\author{Your Name}
\date{2026-01-12}

\begin{document}
\maketitle

{
\setcounter{tocdepth}{2}
\tableofcontents
}
\subsection{🧭 Mission Background}\label{mission-background}

Welcome to your first day as a Junior Analyst for the \textbf{Laser
League}. You've been handed a raw data export of player performance.
Before the Commissioner can use this to rank the world's best players,
you need to prove that the data is structured correctly and free of
``Gremlins.''

\textbf{Your Goal:} Perform a forensic audit. Identify what data types R
assigned, check the table's structural integrity, and document any
quality issues.

\begin{Shaded}
\begin{Highlighting}[]
\CommentTok{\# {-}{-}{-} UTILITY REFERENCE {-}{-}{-}}
\CommentTok{\# library(): Loads the forensic toolkit into your session.}
\CommentTok{\# readr: Part of the Tidyverse, used for fast and "smart" data loading.}
\CommentTok{\# dplyr: Provides the tools for inspection (glimpse) and manipulation.}

\FunctionTok{library}\NormalTok{(readr)}
\FunctionTok{library}\NormalTok{(dplyr)}

\CommentTok{\# Define the root path (Adjust if your folder structure is different)}
\NormalTok{ENIA\_ROOT }\OtherTok{\textless{}{-}} \StringTok{"\textasciitilde{}/ENIA{-}5120"}
\NormalTok{league\_path }\OtherTok{\textless{}{-}} \FunctionTok{file.path}\NormalTok{(}\StringTok{"../../../data"}\NormalTok{, }\StringTok{"laser\_league\_players.csv"}\NormalTok{)}
\end{Highlighting}
\end{Shaded}

\subsection{1. The Blueprint: Loading \&
Inspection}\label{the-blueprint-loading-inspection}

Our first task is to bring the evidence onto our digital workbench (the
R Environment).

\subsubsection{1.1 Load the Evidence}\label{load-the-evidence}

\textbf{functions:}\\
\texttt{read\_csv()}: Loads data much faster and smarter than base R's
read.csv().\\
\texttt{head()}: Displays the first 6 rows of a dataset.

\textbf{logic steps:}\\
Use \texttt{read\_csv()} to load the file located at
\texttt{league\_path}.\\
Store the result in an object named \texttt{players}.\\
Use \texttt{head()} to verify that the table looks correct.

\begin{Shaded}
\begin{Highlighting}[]
\NormalTok{players }\OtherTok{\textless{}{-}} \FunctionTok{read\_csv}\NormalTok{(league\_path)}
\end{Highlighting}
\end{Shaded}

\begin{verbatim}
## Rows: 181 Columns: 21
## -- Column specification --------------------------------------------------------
## Delimiter: ","
## chr (10): player_id, squad_id, squad_name, home_arena, role, preferred_loado...
## dbl (11): matches_played, assists, season_score, tag_accuracy_pct, shield_up...
## 
## i Use `spec()` to retrieve the full column specification for this data.
## i Specify the column types or set `show_col_types = FALSE` to quiet this message.
\end{verbatim}

\begin{Shaded}
\begin{Highlighting}[]
\FunctionTok{head}\NormalTok{(players)}
\end{Highlighting}
\end{Shaded}

\begin{verbatim}
## # A tibble: 6 x 21
##   player_id squad_id squad_name  home_arena   role     preferred_loadout
##   <chr>     <chr>    <chr>       <chr>        <chr>    <chr>            
## 1 LT0001    SQ045    Hyper Hawks astro alley  Scout    Stealth Rig      
## 2 LT0002    SQ045    Hyper Hawks Astro Alley  Captain  Dual Pulse       
## 3 LT0003    SQ041    Neon Ninjas Photon Forge Captain  Support Pod      
## 4 LT0004    SQ041    Neon Ninjas Photon Forge Sniper   Support Pod      
## 5 LT0005    SQ045    Hyper Hawks Nebula Nexus Captain  Rail Blaster     
## 6 LT0006    SQ045    Hyper Hawks ASTRO ALLEY  Engineer Stealth Rig      
## # i 15 more variables: matches_played <dbl>, assists <dbl>, season_score <dbl>,
## #   tag_accuracy_pct <dbl>, shield_uptime_pct <dbl>, stamina_score <dbl>,
## #   reflex_score <dbl>, reaction_time_sec <dbl>, energy_delta_pct <dbl>,
## #   favorite_drink <chr>, preferred_channel <chr>, commute_mode <chr>,
## #   glow_rating <dbl>, wellbeing_index <dbl>, scout_notes <chr>
\end{verbatim}

\subsubsection{1.2 Identify the Data
Structure}\label{identify-the-data-structure}

\textbf{functions:}\\
\texttt{class()}: Tells you the ``Object Type'' (e.g.,
tibble/data.frame).\\
\texttt{glimpse()}: Provides a clean, vertical overview of column names
and variable types.

\textbf{logic steps:}\\
Check the \texttt{class()} of your \texttt{players} object to see if R
treats it as a table.\\
Use \texttt{glimpse()} to see every column name and its assigned data
type (e.g., , ).

\begin{Shaded}
\begin{Highlighting}[]
\FunctionTok{class}\NormalTok{(players) }
\end{Highlighting}
\end{Shaded}

\begin{verbatim}
## [1] "spec_tbl_df" "tbl_df"      "tbl"         "data.frame"
\end{verbatim}

\begin{Shaded}
\begin{Highlighting}[]
\FunctionTok{glimpse}\NormalTok{(players)}
\end{Highlighting}
\end{Shaded}

\begin{verbatim}
## Rows: 181
## Columns: 21
## $ player_id         <chr> "LT0001", "LT0002", "LT0003", "LT0004", "LT0005", "L~
## $ squad_id          <chr> "SQ045", "SQ045", "SQ041", "SQ041", "SQ045", "SQ045"~
## $ squad_name        <chr> "Hyper Hawks", "Hyper Hawks", "Neon Ninjas", "Neon N~
## $ home_arena        <chr> "astro alley", "Astro Alley", "Photon Forge", "Photo~
## $ role              <chr> "Scout", "Captain", "Captain", "Sniper", "Captain", ~
## $ preferred_loadout <chr> "Stealth Rig", "Dual Pulse", "Support Pod", "Support~
## $ matches_played    <dbl> 22, 7, 32, 10, 6, 7, 39, 11, 25, 12, 23, 37, 25, 19,~
## $ assists           <dbl> 6, 5, 34, 23, 16, 35, 21, 34, 5, 3, 7, 17, 24, 18, 2~
## $ season_score      <dbl> 6710.7, 2921.8, 8553.8, 7615.3, 4008.4, 6561.7, 6215~
## $ tag_accuracy_pct  <dbl> 63.5, 28.5, 58.6, 63.2, 38.2, 41.0, 27.5, 39.0, 47.5~
## $ shield_uptime_pct <dbl> 13.2, 77.2, 60.7, 94.8, 33.1, 12.0, 24.2, 85.9, 21.3~
## $ stamina_score     <dbl> 84, 51, 81, 75, 77, 90, 87, 71, 84, 92, 70, 51, 66, ~
## $ reflex_score      <dbl> 75, 58, 73, 74, 80, 50, 64, 45, 48, 94, 43, 67, 58, ~
## $ reaction_time_sec <dbl> 0.264, 0.346, 0.311, 0.250, 0.292, 0.392, 0.336, 0.3~
## $ energy_delta_pct  <dbl> 20.6, 12.2, 6.2, 19.7, -1.7, 7.6, 17.2, 10.2, 9.5, 5~
## $ favorite_drink    <chr> "Nebula Nectar", "Quasar Quench", "Photon Fizz", "Ph~
## $ preferred_channel <chr> "Arena App", "Squad Mentor", "Arena App", "Beacon Sc~
## $ commute_mode      <chr> "Carpool", "Hover Tram", "Hover Tram", "Transit", "C~
## $ glow_rating       <dbl> 7.5, 1.2, 9.9, 9.1, 3.0, 6.5, 6.3, 7.5, 5.4, 3.9, 4.~
## $ wellbeing_index   <dbl> 85.7, 64.1, 77.3, 81.1, 70.2, 84.1, 87.0, 74.7, 81.6~
## $ scout_notes       <chr> "needs recalibration", NA, "late check-in", NA, NA, ~
\end{verbatim}

\textbf{Question 1.2.1:} What R class is players? How does R internally
store this table?

\textbf{Question 1.2.2:} Based on glimpse(), how many rows
(observations) and columns (variables) are present?

\subsection{2. The Language: Data Type
Checks}\label{the-language-data-type-checks}

In R, the type of data determines what math is allowed. You cannot
calculate the ``average'' of a word.

\subsubsection{2.1 Inspect Individual
Columns}\label{inspect-individual-columns}

\textbf{functions:}\\
\texttt{class(df\$column)}: Extracts a specific column using the \$ sign
and reveals its type.

\textbf{logic steps:}\\
Use the \texttt{\$} operator to check the class of
\texttt{player\_id}.\\
Check the class of \texttt{season\_score}.\\
Check the class of \texttt{reaction\_time\_sec}.

\begin{Shaded}
\begin{Highlighting}[]
\FunctionTok{class}\NormalTok{(players}\SpecialCharTok{$}\NormalTok{player\_id)}
\end{Highlighting}
\end{Shaded}

\begin{verbatim}
## [1] "character"
\end{verbatim}

\begin{Shaded}
\begin{Highlighting}[]
\FunctionTok{class}\NormalTok{(players}\SpecialCharTok{$}\NormalTok{season\_score)}
\end{Highlighting}
\end{Shaded}

\begin{verbatim}
## [1] "numeric"
\end{verbatim}

\begin{Shaded}
\begin{Highlighting}[]
\FunctionTok{class}\NormalTok{(players}\SpecialCharTok{$}\NormalTok{reaction\_time\_sec)}
\end{Highlighting}
\end{Shaded}

\begin{verbatim}
## [1] "numeric"
\end{verbatim}

\textbf{Question 2.1.1:} assists shows up as in glimpse(). Is this
Nominal, Ordinal, or Numeric? Why?

\textbf{Question 2.1.2:} role is stored as a character. Given there's no
ranking among roles, is this Nominal or Ordinal? Why does that matter
for our future analysis?

\subsubsection{2.2 Numerical Summaries}\label{numerical-summaries}

\textbf{functions:}\\
\texttt{summary()}: Provides the Min, Max, Median, and Mean.

\textbf{logic steps:}\\
Run a \texttt{summary()} on the \texttt{season\_score} column.\\
Look at the Max value---if it is impossibly high, you've found a
``Gremlin.''

\begin{Shaded}
\begin{Highlighting}[]
\FunctionTok{summary}\NormalTok{(players}\SpecialCharTok{$}\NormalTok{season\_score)}
\end{Highlighting}
\end{Shaded}

\begin{verbatim}
##    Min. 1st Qu.  Median    Mean 3rd Qu.    Max. 
##    2650    4988    5868    5788    6795    8703
\end{verbatim}

\textbf{Question 2.2.1:} What is the median and maximum season\_score?
Could a massive max value indicate a data quality issue or just a very
good player?

\subsection{3. Forensic Audit: Finding the
Gremlins}\label{forensic-audit-finding-the-gremlins}

Now we look for the quality issues we discussed: Missingness,
Duplicates, and Inconsistent Casing.

\subsubsection{3.1 Missing Values (NA)}\label{missing-values-na}

\textbf{functions:}\\
\texttt{is.na()}: Returns TRUE for every missing value.\\
\texttt{colSums()}: Adds up all the TRUEs to give a total count of
missing values per column.

\textbf{logic steps:}\\
Combine \texttt{is.na(players)} with \texttt{colSums()} to get a report
of all missing values in the entire table.

\begin{Shaded}
\begin{Highlighting}[]
\FunctionTok{colSums}\NormalTok{(}\FunctionTok{is.na}\NormalTok{(players))}
\end{Highlighting}
\end{Shaded}

\begin{verbatim}
##         player_id          squad_id        squad_name        home_arena 
##                 0                 0                 0                 0 
##              role preferred_loadout    matches_played           assists 
##                 0                31                 0                 0 
##      season_score  tag_accuracy_pct shield_uptime_pct     stamina_score 
##                 0                 0                 0                 0 
##      reflex_score reaction_time_sec  energy_delta_pct    favorite_drink 
##                 0                 0                 0                 2 
## preferred_channel      commute_mode       glow_rating   wellbeing_index 
##                 0                 0                 0                 0 
##       scout_notes 
##                27
\end{verbatim}

\textbf{Question 3.1.1:} Which variable has the most missing values? If
this were the primary \(CO_2\) measurement in a climate study, why would
that be a problem?

\subsubsection{3.2 Inconsistent Casing}\label{inconsistent-casing}

\textbf{functions:}\\
\texttt{count()}: Groups the rows by a column and counts how many times
each label appears.\\
\texttt{arrange(desc(n))}: Sorts the counts from highest to lowest.

\textbf{logic steps:}\\
Take the \texttt{players} data.\\
AND THEN (\%\textgreater\%) \texttt{count()} the \texttt{home\_arena}
column.\\
AND THEN (\%\textgreater\%) \texttt{arrange()} by descending count
(\texttt{n}).

\begin{Shaded}
\begin{Highlighting}[]
\NormalTok{players }\SpecialCharTok{\%\textgreater{}\%}
  \FunctionTok{count}\NormalTok{(home\_arena) }\SpecialCharTok{\%\textgreater{}\%}
  \FunctionTok{arrange}\NormalTok{(}\FunctionTok{desc}\NormalTok{(n))}
\end{Highlighting}
\end{Shaded}

\begin{verbatim}
## # A tibble: 8 x 2
##   home_arena       n
##   <chr>        <int>
## 1 Astro Alley     51
## 2 Nebula Nexus    47
## 3 Photon Forge    42
## 4 Nova Spire      37
## 5 ASTRO ALLEY      1
## 6 Photon forge     1
## 7 astro alley      1
## 8 photon Forge     1
\end{verbatim}

\textbf{Question 3.2.1:} Do you notice inconsistent casing or spelling
(e.g., neon yard vs Neon Yard)? How would this mess up a bar chart of
arena popularity?

\subsubsection{3.3 Duplicate IDs}\label{duplicate-ids}

\textbf{functions:}\\
\texttt{duplicated()}: Flags rows that have appeared before in the
dataset.

\textbf{logic steps:}\\
Use \texttt{sum(duplicated())} on the \texttt{player\_id} column to see
if any unique identifier appears more than once.

\begin{Shaded}
\begin{Highlighting}[]
\FunctionTok{sum}\NormalTok{(}\FunctionTok{duplicated}\NormalTok{(players}\SpecialCharTok{$}\NormalTok{player\_id))}
\end{Highlighting}
\end{Shaded}

\begin{verbatim}
## [1] 1
\end{verbatim}

\subsection{4. Challenge: Spot the
Issues}\label{challenge-spot-the-issues}

Look at this slice of ``Raw Data'' from the League's backup server.

Plaintext

\begin{longtable}[]{@{}llll@{}}
\toprule\noalign{}
player\_id & home\_arena & season\_score & preferred\_loadout \\
\midrule\noalign{}
\endhead
\bottomrule\noalign{}
\endlastfoot
P-0220 & neon yard & 905 & Hyper Beam \\
P-0455 & Glow Hall & NA & Pulse Shield \\
P-0220 & Neon Yard & 905 & Hyper Beam \\
P-0512 & Photon Dome & -10 & '' '' \\
\end{longtable}

\textbf{Question 4.1.1:} List three distinct data quality issues from
this snippet (e.g., Impossible numbers, Duplicates, Casing).

\textbf{Question 4.1.2:} Which specific R functions from this lab would
you use to confirm each issue if it appeared in a massive dataset of
100,000 players?

\end{document}
